\subsubsection{Data Management}
\begin{enumerate}
    \item \textbf{Data Reading and Processing}
          \begin{itemize}
            \item \textbf{Read Gene IDs from TSV File}
                  
                  The function reads gene identifiers from a tab-separated value (TSV) file,
                  extracting specific gene IDs from a designated column while skipping
                  header rows. This function is typically used as the first step in processing
                  gene expression data to obtain the complete list of gene identifiers.
                  Call the \texttt{read\_gene\_ids\_from\_tsv\_file} function with the following
                  arguments:
                  
                  \begin{itemize}
                    \item \texttt{filename}: Path to the TSV file containing gene IDs (string
                          or list with single file)
                          
                    \item \texttt{n\_genes}: Total number of genes in the file
                          
                    \item \texttt{gene\_ids\_len}: Maximum character length for each gene
                          ID
                          
                    \item \texttt{n\_header\_rows}: Number of header rows to skip at the
                          beginning of the file
                          
                    \item \texttt{gene\_col}: Column index (1-based) where gene IDs are located
                  \end{itemize}
                  
                  \textbf{Returns:}
                  \begin{itemize}
                    \item \textbf{gene\_ids}: Numpy array of shape \texttt{(n\_genes)}
                          containing gene ID strings
                  \end{itemize}
                  
                  \begin{lstlisting}[language=Python, caption={Python Read Gene IDs from TSV File}]
gene_ids = read_gene_ids_from_tsv_file(filename, n_genes,
                            gene_ids_len, n_header_rows, gene_col)
                  \end{lstlisting}
                  
            \item \textbf{Read Expression Vectors}
                  
                  The function reads gene expression data from multiple tabular files (CSV/TSV
                  format), extracting expression values for specified genes across
                  multiple samples. This function consolidates expression data from multiple
                  files into a single 2D array. Call the \texttt{read\_expression\_vectors}
                  function with the following arguments:
                  
                  \begin{itemize}
                    \item \texttt{file\_list}: List of file paths containing expression data
                          
                    \item \texttt{gene\_ids}: Array of gene IDs to extract expression values
                          for
                          
                    \item \texttt{n\_samples}: Number of samples
                          
                    \item \texttt{n\_header\_rows}: Number of header rows to skip in each
                          file
                          
                    \item \texttt{gene\_col}: Column index (1-based) containing gene identifiers
                          
                    \item \texttt{value\_cols}: List of column indices (1-based) containing
                          expression values
                          
                    \item \textbf{(Optional)} \texttt{delimiter}: Delimiter used in the
                          files (default: tab)
                  \end{itemize}
                  
                  \textbf{Returns:}
                  \begin{itemize}
                    \item \textbf{expression\_vectors}: 2D numpy array of shape \texttt{(n\_samples,
                          n\_genes)} containing expression values as 64-bit floats
                  \end{itemize}
                  
                  \begin{lstlisting}[language=Python, caption={Python Read Expression Vectors}]
expression_vectors = read_expression_vectors(file_list, gene_ids, 
                                n_samples, n_header_rows, gene_col, 
                                value_cols, delimiter='\t')
                  \end{lstlisting}
                  
            \item \textbf{Read OrthoFinder File}
                  
                  The function processes an OrthoFinder output file to map genes to their
                  respective orthologous families. This function establishes the gene-family
                  relationships essential for downstream orthology-based analyses. Call the
                  \texttt{read\_orthofinder\_file} function with the following arguments:
                  
                  \begin{itemize}
                    \item \texttt{filename}: Path to the OrthoFinder TSV file (string or
                          list with single file)
                          
                    \item \texttt{gene\_ids}: Complete list of gene identifiers to map
                          
                    \item \texttt{family\_ids\_len}: Maximum character length for family
                          IDs
                          
                    \item \texttt{n\_families}: Total number of gene families expected
                  \end{itemize}
                  
                  \textbf{Return dictionary:}
                  \begin{itemize}
                    \item \textbf{family\_ids}: Numpy array of shape \texttt{(n\_families,)}
                          containing family ID strings
                          
                    \item \textbf{gene\_to\_fam}: Numpy array of shape \texttt{(n\_genes,)}
                          containing family indices for each gene (0 = unassigned)
                  \end{itemize}
                  
                  \begin{lstlisting}[language=Python, caption={Python Read OrthoFinder File}]
family_result = read_orthofinder_file(filename, gene_ids, 
                                family_ids_len, n_families)
                  \end{lstlisting}
                  
                  \textbf{Complete Data Loading Workflow Example:} \begin{lstlisting}[language=Python, caption={Python Complete Data Loading Workflow}]
# Step 1: Load gene identifiers
gene_ids = read_gene_ids_from_tsv_file(
    filename="data/expression_data.tsv",
    n_genes=50000,
    gene_ids_len=32,
    n_header_rows=1,
    gene_col=1
)

# Step 2: Load expression data from multiple samples
samples = ["sample1.tsv", "sample2.tsv", "sample3.tsv", "sample4.tsv"]
expression_vectors = read_expression_vectors(
    file_list=samples,
    gene_ids=gene_ids,
    n_samples=12, # assuming 3 samples per file
    n_header_rows=1,
    gene_col=1,
    value_cols=[2, 3]
)

# Step 3: Load orthology information
orthology_data = read_orthofinder_file(
    filename="data/Orthogroups.tsv",
    gene_ids=gene_ids,
    family_ids_len=32,
    n_families=15000
)
            \end{lstlisting}
          \end{itemize}
          
    \item \textbf{Data Filtering and Validation}
          \begin{itemize}
            \item \textbf{Filter Unassigned Genes}
                  
                  The function filters out genes that are not assigned to any family (where
                  gene\_to\_fam = 0). This is essential for cleaning data before
                  orthology-based analyses, ensuring only genes with valid family assignments
                  are processed. Call the \texttt{filter\_unassigned\_genes} function with
                  the following arguments:
                  
                  \begin{itemize}
                    \item \texttt{gene\_to\_fam}: Family assignment array where 0 indicates
                          unassigned genes
                  \end{itemize}
                  
                  \textbf{Return dictionary:}
                  \begin{itemize}
                    \item \textbf{mask}: Boolean array indicating which genes to keep (1
                          = assigned to family, 0 = unassigned)
                          
                    \item \textbf{n\_genes\_kept}: Number of genes kept after filtering
                  \end{itemize}
                  
                  \begin{lstlisting}[language=Python, caption={Python Filter Unassigned Genes}]
filtered_genes = filter_unassigned_genes(gene_to_fam)
                  \end{lstlisting}
                  
            \item \textbf{Validate Gene to Family Mapping}
                  
                  The function validates the gene-to-family mapping array, ensuring all
                  family indices are within valid range and properly assigned. This prevents
                  downstream errors from invalid family indices. Call the \texttt{validate\_gene\_to\_family\_mapping}
                  function with the following arguments:
                  
                  \begin{itemize}
                    \item \texttt{gene\_to\_fam}: Array mapping each gene to a family index
                          (0 = unassigned)
                          
                    \item \texttt{n\_families}: Total number of families
                  \end{itemize}
                  
                  \begin{lstlisting}[language=Python, caption={Python Validate Gene to Family Mapping}]
validate_gene_to_family_mapping(gene_to_fam, n_families)
                  \end{lstlisting}
                  
            \item \textbf{Validate Expression Data}
                  
                  The function validates expression data, checking for NaN/Inf values
                  and optionally verifying non-negative values. This ensures data quality
                  before computational analyses. Call the \texttt{validate\_expression\_data}
                  function with the following arguments:
                  
                  \begin{itemize}
                    \item \texttt{expression\_vectors}: 2D array of expression data (n\_sampels
                          x n\_genes)
                          
                    \item \textbf{(Optional)} \texttt{check\_non\_negative}: Whether to
                          check for non-negative values (default: True)
                  \end{itemize}
                  
                  \begin{lstlisting}[language=Python, caption={Python Validate Expression Data}]
validate_expression_data(expression_vectors, check_non_negative=True)
                  \end{lstlisting}
                  
            \item \textbf{Validate Family Centroids}
                  
                  The function validates family centroids data, checking for NaN/Inf
                  values and proper data structure. Centroids represent the central tendency
                  of each gene family's expression profile. Call the \texttt{validate\_family\_centroids}
                  function with the following arguments:
                  
                  \begin{itemize}
                    \item \texttt{family\_centroids}: 2D array of family centroids (samples
                          x families)
                  \end{itemize}
                  
                  \begin{lstlisting}[language=Python, caption={Python Validate Family Centroids}]
validate_family_centroids(family_centroids)
                  \end{lstlisting}
                  
            \item \textbf{Validate Shift Vectors}
                  
                  The function validates shift vectors, ensuring data types are correct
                  and structure matches expected patterns. Shift vectors represent the deviation
                  of each gene from its family centroid. Call the \texttt{validate\_shift\_vectors}
                  function with the following arguments:
                  
                  \begin{itemize}
                    \item \texttt{shift\_vectors}: 2D array of shift vectors (samples ×
                          genes)
                          
                    \item \texttt{expression\_vectors}: 2D array of expression data (samples
                          × genes)
                          
                    \item \texttt{family\_centroids}: 2D array of family centroids (samples
                          × families)
                          
                    \item \texttt{gene\_to\_fam}: Array mapping each gene to a family index
                          
                    \item \texttt{n\_genes}: Number of genes
                          
                    \item \texttt{n\_samples}: Number of samples
                          
                    \item \texttt{n\_families}: Number of families
                  \end{itemize}
                  
                  \begin{lstlisting}[language=Python, caption={Python Validate Shift Vectors}]
validate_shift_vectors(shift_vectors, expression_vectors, 
                    family_centroids, gene_to_fam, n_genes, 
                    n_samples, n_families)
                  \end{lstlisting}
                  
            \item \textbf{Validate String Array Uniqueness}
                  
                  The function validates that all strings in an array are unique, using
                  a hashset internally. Note: This may increase memory usage temporarily
                  for large datasets. Call the \texttt{validate\_string\_array\_uniqueness}
                  function with the following arguments:
                  
                  \begin{itemize}
                    \item \texttt{strings}: Array of strings to validate (typically gene
                          IDs or family IDs)
                  \end{itemize}
                  
                  \begin{lstlisting}[language=Python, caption={Python Validate String Array Uniqueness}]
validate_string_array_uniqueness(strings)
                  \end{lstlisting}
                  
            \item \textbf{Validate Data Structure}
                  
                  The function performs comprehensive validation of the overall data structure
                  consistency, confirming sizes and dependencies between different data
                  components as far as possible. Call the \texttt{validate\_data\_structure}
                  function with the following arguments:
                  
                  \begin{itemize}
                    \item \texttt{n\_genes}: Number of genes
                          
                    \item \texttt{n\_families}: Number of families
                          
                    \item \texttt{n\_samples}: Number of samples
                          
                    \item \texttt{gene\_ids}: Array of gene IDs
                          
                    \item \texttt{gene\_family\_ids}: Array of family IDs
                          
                    \item \texttt{gene\_to\_fam}: Array mapping each gene to a family index
                          
                    \item \texttt{expression\_vectors}: 2D array of expression data (samples
                          × genes)
                          
                    \item \texttt{family\_centroids}: 2D array of family centroids (samples
                          × families)
                          
                    \item \texttt{shift\_vectors}: 2D array of shift vectors (genes ×
                          samples)
                  \end{itemize}
                  
                  \begin{lstlisting}[language=Python, caption={Python Validate Data Structure}]
validate_data_structure(n_genes, n_families, n_samples, gene_ids, 
                gene_family_ids, gene_to_fam, expression_vectors, 
                family_centroids, shift_vectors)
                  \end{lstlisting}
                  
            \item \textbf{Validate All Data}
                  
                  The function performs comprehensive validation of all data components in
                  one call. This function executes all individual validations
                  sequentially for complete data quality assurance. Call the \texttt{validate\_all\_data}
                  function with the following arguments:
                  
                  \begin{itemize}
                    \item \texttt{n\_genes}: Number of genes
                          
                    \item \texttt{n\_families}: Number of families
                          
                    \item \texttt{n\_samples}: Number of samples
                          
                    \item \texttt{gene\_ids}: Array of gene IDs
                          
                    \item \texttt{gene\_family\_ids}: Array of family IDs
                          
                    \item \texttt{gene\_to\_fam}: Array mapping each gene to a family index
                          
                    \item \texttt{expression\_vectors}: 2D array of expression data (samples
                          × genes)
                          
                    \item \texttt{family\_centroids}: 2D array of family centroids (samples
                          × families)
                          
                    \item \texttt{shift\_vectors}: 2D array of shift vectors (samples ×
                          genes)
                  \end{itemize}
                  
                  \begin{lstlisting}[language=Python, caption={Python Validate All Data}]
validate_all_data(n_genes, n_families, n_samples, gene_ids, 
                    gene_family_ids, gene_to_fam, expression_vectors, 
                    family_centroids, shift_vectors)
                  \end{lstlisting}
                  
                  \textbf{Data Filtering and Validation Workflow Example:} 
                  \begin{lstlisting}[language=Python, caption={Python Data Filtering and Validation Workflow}]
# After loading data, filter unassigned genes
filter_result = filter_unassigned_genes(gene_to_fam)
mask = np.array(filter_result['mask'], dtype=bool)
n_genes_kept = filter_result['n_genes_kept']

# Apply filter to all gene-based arrays
filtered_gene_ids = gene_ids[mask]
filtered_expression = expression_vectors[:, mask]
filtered_gene_to_fam = gene_to_fam[mask]

# Run comprehensive validations
validate_string_array_uniqueness(filtered_gene_ids)
validate_string_array_uniqueness(family_ids)
validate_expression_data(filtered_expression, check_non_negative=True)
validate_gene_to_family_mapping(filtered_gene_to_fam, n_families)

# For complete analysis with centroids and shift vectors
validate_data_structure(
    n_genes_kept, n_families, n_samples, 
    filtered_gene_ids, family_ids, filtered_gene_to_fam,
    filtered_expression, family_centroids, shift_vectors
)

# Or use the comprehensive validation
validate_all_data(
    n_genes_kept, n_families, n_samples,
    filtered_gene_ids, family_ids, filtered_gene_to_fam,
    filtered_expression, family_centroids, shift_vectors
)
            \end{lstlisting}
          \end{itemize}
\end{enumerate}